% ======================================================================
%  A Ground Truth Dataset for Evaluating Agentic Attack Surface
%  Enumeration Across Vulnerable Applications
%  FYP research report -- YEO Dominic (SMU SCIS)
%
%  Compile:  pdflatex gt_report.tex  (run twice for refs/ToC)
%  Self-contained: no external image files required.
% ======================================================================
\documentclass[11pt,a4paper]{article}

\usepackage[T1]{fontenc}
\usepackage[utf8]{inputenc}
\usepackage[margin=2.4cm]{geometry}
\usepackage{microtype}
\usepackage{booktabs}
\usepackage{array}
\usepackage{longtable}
\usepackage{tabularx}
\usepackage{amsmath,amssymb}
\usepackage{enumitem}
\usepackage{xcolor}
\usepackage{listings}
\usepackage{caption}
\usepackage{graphicx}
\usepackage{tikz}
\usetikzlibrary{positioning,arrows.meta,shapes.geometric}
\usepackage[hidelinks]{hyperref}

% ---- palette (deck-inspired navy + gold) -----------------------------
\definecolor{navy}{HTML}{16256A}
\definecolor{gold}{HTML}{B5861C}
\definecolor{inkgray}{HTML}{4A5468}
\definecolor{codebg}{HTML}{F5F6F9}
\definecolor{l1blue}{HTML}{EAF0FB}
\definecolor{l2green}{HTML}{EAF5EE}
\definecolor{agamber}{HTML}{FBF3E1}
\definecolor{redbox}{HTML}{FBEEEB}

\captionsetup{font=small,labelfont={bf,color=navy}}
\setlength{\parskip}{4pt}
\setlength{\parindent}{0pt}

% ---- code listings (verbatim: no escaping needed inside) -------------
\lstdefinestyle{gt}{
  basicstyle=\ttfamily\footnotesize,
  backgroundcolor=\color{codebg},
  frame=single, rulecolor=\color{black!15},
  breaklines=true, columns=fullflexible,
  keepspaces=true, showstringspaces=false,
  xleftmargin=6pt, xrightmargin=6pt, aboveskip=6pt, belowskip=6pt,
}
\lstset{style=gt}
% inline code token helper (handles underscores)
\newcommand{\cd}[1]{\lstinline[basicstyle=\ttfamily\small]|#1|}

\newcommand{\L}[1]{\textbf{L#1}}

% ======================================================================
\begin{document}

\begin{titlepage}
\centering
\vspace*{1.5cm}
{\color{gold}\small\bfseries SINGAPORE MANAGEMENT UNIVERSITY}\\
{\color{inkgray}\small School of Computing and Information Systems}\\[2.4cm]
{\color{navy}\LARGE\bfseries A Ground Truth Dataset for Evaluating\\[4pt]
Agentic Attack Surface Enumeration\\[4pt] Across Vulnerable Applications}\\[1.6cm]
{\large Final Year Project Report}\\[2.2cm]
{\large\bfseries YEO Dominic}\\[6pt]
{\normalsize Supervised by Dr.~XIE Xiaofei and YANG Ruozhao}\\[1.6cm]
{\normalsize November 2026}\\
\vfill
{\footnotesize SMU Classification: Restricted}
\end{titlepage}

% ----------------------------------------------------------------------
\begin{abstract}
\noindent
AI agents are increasingly used to analyse software and explore web applications,
yet existing agent-security benchmarks measure downstream task or vulnerability
outcomes and provide little evidence of how effectively an agent \emph{discovers and
describes} an application's attack surface. This project establishes a ground-truth
foundation for that evaluation. We define a three-layer schema that represents an
application's attack surface as (i)~deployed HTTP interfaces and their request/response
contracts, (ii)~semantic business actions and the state they depend on, and
(iii)~agent capabilities and data lineage. Using a repeatable workflow that reconciles
white-box source analysis with black-box runtime capture and enforces cross-layer
referential integrity, we produce an evidence-backed dataset across \textbf{fifteen}
open-source vulnerable applications (ten traditional web applications and five
LLM-integrated systems). The resulting dataset records \textbf{1{,}397 deployed
endpoints}, \textbf{980 business semantic flows}, and \textbf{327 agentic capability
records}. The dataset is \emph{evidence-graded}: every interface and action carries a
runtime-observation depth and every agentic capability a runtime status, so that
across the traditional-web corpus all 1{,}397 interfaces and 980 actions are
runtime-grounded (the actions split 276 confirmed by observed effect and 704 by
interface), while the agentic layer is grounded for 129 of 307 records with each
remaining record labelled by the precise reason it cannot be observed in a local,
credential-free environment. Beyond the artefact, the project contributes a reusable
annotation-and-review methodology and surfaces two findings: the same business action is frequently exposed
through multiple surfaces that automated crawlers under-count, and the root cause of
agentic attack vectors is \emph{authority}---an untrusted input reaching a privileged
tool---which the schema makes directly enumerable and checkable for enforced controls.
\end{abstract}

\tableofcontents
\clearpage

% ======================================================================
\section{Introduction and Motivation}
% ======================================================================

\subsection{What is an attack surface?}
An application exposes functionality through UI pages, APIs, and state-dependent
workflows. Each of these is an entry point that a specific attack vector can target,
and together they form the application's \emph{attack surface}. LLM-driven applications
expand this surface with non-deterministic interfaces and context pipelines---system
prompts, retrieval-augmented-generation (RAG) datastores, and tool-calling
hooks---that a conventional, route-centric surface map does not capture.

\subsection{AI in attack surface generation}
AI is increasingly used in security tooling to analyse software and explore web
applications. Multimodal language-model crawlers can use page state and functional
context to guide discovery. However, state-dependent workflows, event-driven
interfaces, and API-only functionality that never appears in the browser are routinely
missed by automated discovery.

\subsection{Limitation of existing benchmarks}
Existing agent-security benchmarks tend to measure \emph{downstream} task completion or
vulnerability-exploitation outcomes. They therefore leave limited evidence about how
effectively an agent discovers and describes an application's attack surface---and
discovery recall varies substantially across models, which is precisely the quantity
that a reliable reference is needed to score.

\subsection{Motivation}
This project aims to establish a ground-truth foundation for the future evaluation of
agents' effectiveness at discovering attack surfaces. It gives future evaluations a
reference for assessing not only \emph{which} surfaces an agent finds, but also
\emph{how accurately} it describes their application context and the relationships
between them.

% ======================================================================
\section{Research Scope and Contributions}
% ======================================================================

\subsection{Scope}
\begin{enumerate}[leftmargin=1.6em,itemsep=2pt]
\item Evidence-grounded attack-surface analysis across \textbf{15} open-source
      vulnerable applications---\textbf{10} traditional web applications and \textbf{5}
      LLM-integrated systems---chosen to represent a variety of surface types.
\item The attack surface is represented at three levels:
      (a)~deployed web interactions and their request/response contracts;
      (b)~semantic business actions and the state they depend on;
      (c)~agent capabilities and data lineage.
\item The surface is generated using a combination of black-box dynamic analysis and
      white-box static source analysis, reconciled against one another.
\end{enumerate}

\subsection{Contributions}
\begin{enumerate}[leftmargin=1.6em,itemsep=2pt]
\item \textbf{An evidence-backed, multi-application dataset} recording
      application-specific surfaces with version and deployment context, observed
      behaviour, semantic annotations, and supporting evidence:
      \begin{itemize}[leftmargin=1.4em,itemsep=0pt]
        \item \textbf{1{,}397} deployed endpoints (\L1);
        \item \textbf{980} business semantic flows (\L2);
        \item \textbf{327} agentic capability records.
      \end{itemize}
\item \textbf{A repeatable annotation and review workflow} that reconciles source and
      runtime evidence, generates structured records, tracks unresolved cases, and
      automatically checks cross-layer references.
\item \textbf{An evidence-grading method for runtime grounding.} Each record is assigned
      a runtime-observation depth from a fixed ladder; a completeness oracle measures
      authored surface against each application's own route list; a multi-identity probe
      records per-role reachability; and every agentic record that cannot be observed in
      a local, credential-free environment is tagged with a blocker from a closed
      taxonomy. The grading is honest by construction: a record is marked observed only
      when a real response or effect is captured, never for a source-only path.
\end{enumerate}

% ======================================================================
\section{The Three-Layer Schema}
% ======================================================================
The schema is the central contribution. It is deliberately grounded in recognised
standards so that labels are defensible rather than ad hoc. Figure~\ref{fig:layers}
shows the three layers and the identifier fields that bind them.

\begin{figure}[ht]\centering
\begin{tikzpicture}[
  layer/.style={draw, rounded corners, text width=7.4cm, minimum height=1.2cm,
                align=left, inner sep=6pt, font=\small},
  lbl/.style={font=\scriptsize\itshape, text=inkgray}]
\node[layer, fill=l1blue] (l1)
  {\textbf{\color{navy}L1 --- Deployed HTTP Interfaces}\\[1pt]
   \scriptsize concrete routes, pages, static; \cd{key}, \cd{methods}, \cd{io_variants}, visibility};
\node[layer, fill=l2green, below=1.4cm of l1] (l2)
  {\textbf{\color{navy}L2 --- Semantic Business Actions}\\[1pt]
   \scriptsize Subject--Operation--Object; \cd{requirements}, \cd{effects}, \cd{l1_bindings}};
\node[layer, fill=agamber, below=1.4cm of l2] (ag)
  {\textbf{\color{navy}Agentic --- Capabilities \& Data Lineage}\\[1pt]
   \scriptsize agents, tool\_bindings, context/memory, sinks, edges, controls};
\draw[-{Latex[length=2mm]}, thick, navy] (l1) -- node[right,lbl]{operation\_semantic $\cdot$ l1\_bindings} (l2);
\draw[-{Latex[length=2mm]}, thick, navy] (l2) -- node[right,lbl]{related\_l2\_ids} (ag);
\node[draw, rounded corners, fill=yellow!10, text width=4.1cm, right=1.1cm of l2,
      inner sep=5pt, font=\scriptsize, align=left] (x)
  {\textbf{\color{navy}Cross-layer references}\\[1pt]
   invocation\_surfaces[].l1\_id\\ model\_config.writers[].l1\_id\\
   tool\_bindings.l1\_ids\\ related\_l2\_ids / related\_l2\_steps\\[2pt]
   \textit{declared on agentic records;\\ resolve to real L1/L2 nodes}};
\draw[-{Latex[length=1.6mm]}, dashed, gold] (x.north) to[out=90,in=0] (l1.east);
\draw[-{Latex[length=1.6mm]}, dashed, gold] (x.south) to[out=-90,in=0] (ag.east);
\end{tikzpicture}
\caption{The three-layer schema. L1 records deployed interfaces; L2 the business
actions that realise them (bound back to L1 via \cd{l1_bindings}); and the agentic
layer the capabilities and data lineage. A small set of identifier fields stitches the
layers into one connected, validatable graph.}
\label{fig:layers}
\end{figure}

\subsection{L1 --- deployed interface surface}
\L1 describes deployed interfaces and their observable interactions---concrete routes,
pages, and static resources. It maps \emph{where} an actor can interact with the
deployment, \emph{what} can be sent or received, and \emph{under what conditions} each
interaction point is reachable.

\paragraph{How it is built.} Fields are derived from OpenAPI/Swagger and REST/HTTP
semantics. The security labels draw on OWASP: object-tracking labels (\cd{mergeable},
\cd{reference_capable}, \cd{value_kind}$\in$\{identifier, attribute, \dots\}) follow the
OWASP API Security Top~10 entry API1:2023 \emph{Broken Object Level Authorization}; the
visibility/reachability labels follow OWASP WSTG authentication/authorization
state modelling.

\begin{table}[ht]\centering\small
\caption{L1 field groups.}\label{tab:l1}
\begin{tabularx}{\textwidth}{@{}l>{\raggedright\arraybackslash}X@{}}
\toprule
\textbf{Group} & \textbf{Fields (selected)}\\
\midrule
Node identity   & \cd{key} (normalized URL template), \cd{kind}$\in$\{api,page,static\}, \cd{operation_semantic}\\
HTTP surface    & \cd{methods}, \cd{statuses} (runtime-observed), \cd{content_types}, \cd{params}\\
Method IO       & per \cd{io_variant}: \cd{method}, \cd{request}\{\cd{body_format}, \cd{body_example}, \cd{fields[]}\}, \cd{response}\{\cd{format}, value fields, \cd{fields[]}\}, \cd{reason}/\cd{evidence}\\
Request field   & \cd{name}, \cd{location}$\in$\{path,query,header,body\}, \cd{type}, \cd{required}, \cd{flow}, \cd{description}\\
Semantic IO     & \cd{semantic_inputs}/\cd{semantic_outputs}; \cd{semantic_*_decisions[]}\{\cd{raw_field}, \cd{resource}, \cd{mergeable}, \cd{value_kind}, \cd{reference_capable}, \cd{reference_fields}\}\\
Runtime/identity& \cd{samples[]}, \cd{identity_statuses} (role$\to$statuses), \cd{sources}, \cd{referrers}, \cd{resource_types}, \cd{js_sources}, \cd{runtime_api_prefixes}\\
Visibility      & \cd{gt.visibility}\{\cd{minimum}$\in$\{black\_box,grey\_box,white\_box\}\}, \cd{gt.black_box_reachable}, \cd{gt.reachability_conditions}\\
Evidence/notes  & \cd{gt.evidence[]}\{\cd{type},\cd{ref},\cd{supports}\}, \cd{gt.notes}, \cd{edges[]}, \cd{excluded_items[]}\\
\bottomrule
\end{tabularx}
\end{table}

\subsection{L2 --- semantic business actions}
\L2 describes business actions and the states they depend on; one action may be realised
by several \L1 calls. It maps what an actor can cause the application to do, what objects
they can act on, and how the action depends on prior state.

\paragraph{How it is built.} The descriptor structure is drawn from access-control /
authorization modelling (Subject--Operation--Object). The requirements/effects
structure---used to chain operations into multi-step semantic flows---is borrowed from
classical AI planning (STRIPS/PDDL). Facts use a closed predicate vocabulary
(\cd{Exists}, \cd{Know}, \cd{State}, \cd{Relation}, \cd{Can}, \cd{Authenticated}).

\begin{table}[ht]\centering\small
\caption{L2 field groups.}\label{tab:l2}
\begin{tabularx}{\textwidth}{@{}l>{\raggedright\arraybackslash}X@{}}
\toprule
\textbf{Group} & \textbf{Fields (selected)}\\
\midrule
Subject       & \cd{id}, \cd{identity}, \cd{role}, \cd{auth_context} --- distinct subjects are distinct actions\\
Operation     & \cd{business} verb, \cd{crud}$\in$\{Create,Read,Update,Delete\}\\
Object        & \cd{kind}$\in$\{entity,state,relation\}, \cd{type}, \cd{role}; \cd{binding}\{\cd{source},\cd{field},\cd{resource}\}; \cd{participants}, \cd{constraints}\\
Requirements  & \cd{requirements[]}\{\cd{type}$\in$\{authentication,knowledge,existence,\dots\}, \cd{fact}, \cd{reason}\}\\
Effects       & \cd{effects.gen[]}\{\cd{fact}\{\cd{predicate},\cd{arguments}\}, \cd{basis}\}; \cd{effects.kill[]} (old state on Update, Exists on Delete)\\
IO accounting & \cd{io_accounting[]}\{\cd{direction}, \cd{field}, \cd{disposition}$\in$\{req,object\_attribute,gen,observed\_only\}, \cd{facts}\}\\
L1 binding    & \cd{l1_bindings[]}\{\cd{l1_key}, \cd{methods}, \cd{selector}\{\cd{field},\cd{values}\}\}\\
Visibility    & \cd{gt.visibility}, \cd{gt.black_box_reachable}, \cd{gt.reachability_conditions}, \cd{gt.observable_evidence[]}\\
\bottomrule
\end{tabularx}
\end{table}

\subsection{Agentic layer --- capabilities and data lineage}
The agentic layer describes agent capabilities and data paths. It captures how external
or untrusted content can influence a decision, how authority is applied to tool calls,
and how data or side effects move across multiple steps.

\paragraph{How it is built.} The layer is derived from the OWASP Top~10 for LLM
Applications (2025); the structure for tracing an object from entry to exit is derived
from MITRE ATLAS.

\begin{table}[ht]\centering\small
\caption{Agentic record types and their principal fields.}\label{tab:ag}
\begin{tabularx}{\textwidth}{@{}l>{\raggedright\arraybackslash}X@{}}
\toprule
\textbf{Record} & \textbf{Fields (selected)}\\
\midrule
Agents          & identity (\cd{agent_kind}, \cd{purpose}); \cd{invocation_surfaces}, \cd{model_config}, \cd{instruction_sources}; wiring (\cd{tool_binding_ids}, \cd{context_source_ids}, \cd{memory_ids}, \cd{output_sink_ids}, \cd{control_ids}); \cd{resource_limits}, \cd{authority_context}\\
Tool bindings   & \cd{tool_name}, \cd{tool_source}, \cd{server_or_provider}, \cd{schema_ref}; \cd{description_source}, \cd{capabilities}; \cd{execution_identity}, \cd{argument_handling}, \cd{effect_class}, \cd{execution_environment}, \cd{approval_required}; \cd{output_ids}, \cd{control_ids}\\
Context sources & \cd{source_kind}, \cd{source_ref}, \cd{writers}, \cd{provenance}; \cd{destination_agent_ids}, \cd{trust_class}, \cd{transformation}; \cd{persistence_scope}\\
Memory stores   & \cd{store_kind}, \cd{owner_component}; \cd{readers}, \cd{writers}, \cd{write_provenance}, \cd{read_path}; \cd{scope}, \cd{isolation_enforcement}, \cd{retention_and_deletion}\\
Output sinks    & \cd{sink_kind}, \cd{producer_ids}, \cd{audience}; \cd{delivery}, \cd{transformation}, \cd{render_side_effects}; \cd{becomes_context_source_ids}\\
Execution edges & \cd{source}, \cd{destination}; \cd{decision_maker}, \cd{condition}; \cd{input_source_ids}, \cd{argument_sources}, \cd{payload}, \cd{output_ids}; \cd{control_ids}, \cd{validation_or_approval}, \cd{execution_mode}, \cd{loop_bound}\\
Controls        & \cd{control_kind}; \cd{applies_at}, \cd{enforcement_mode}, \cd{implementation}, \cd{failure_mode}; \cd{applies_to_agent_ids}, \cd{applies_to_tool_binding_ids}\\
\bottomrule
\end{tabularx}
\end{table}

% ======================================================================
\section{Methodology}
% ======================================================================
The dataset is produced in four stages---define the schema, choose the environments,
generate the surfaces, and validate---operationalised as two per-application pipelines.

\subsection{Traditional-web generation pipeline}
\begin{enumerate}[leftmargin=1.7em,itemsep=2pt]
\item \textbf{Freeze and deploy.} Pin the exact version, stand up a local instance, and
      record it in \cd{environment.yaml}; the attack surface is tied to that commit.
\item \textbf{White-box enumeration.} Read the server handlers, OpenAPI spec, and
      frontend, and manually list every interface as a curated operation table in code.
\item \textbf{Black-box capture.} Drive the live application as anonymous, user, and
      admin identities, recording traffic as HAR files; normalise concrete paths
      (\cd{/123}$\to$\cd{/{n}}) and aggregate observed statuses, content-types,
      parameter names, and resource types per (method, path).
\item \textbf{Build L1.} Join the source operation table with the HAR verification to
      complete each \L1 row.
\item \textbf{Introduce semantics.} Lift every \cd{flow=true} field to a semantic input
      or output (e.g.\ a token field becomes a reference-capable identifier resource).
\item \textbf{Generate L2.} Map the semantic labels onto a business action
      (Subject--Operation--Object with requirements and effects).
\item \textbf{Validate} (Section~\ref{sec:val}).
\end{enumerate}

\begin{figure}[ht]\centering\footnotesize
\begin{tabular}{@{}p{3.3cm}p{3.6cm}p{3.6cm}p{3.6cm}@{}}
\toprule
\textbf{Evidence} & \textbf{L1 node} & \textbf{Semantics (flow=true)} & \textbf{L2 action}\\
\midrule
\ttfamily GET /api/v1/media/\allowbreak Yt7r1RFGP $\to$ 200 \{friendly\_token, title, state, views\}
&
\ttfamily key: /api/v1/media/\{friendly\_token\}; kind: api; methods:[GET]; io\_variants: req\{friendly\_token: path, flow:true\}
&
\ttfamily friendly\_token $\to$ MediaToken(media) [identifier, reference\_capable]; title $\to$ MediaTitle(media) [attribute]
&
\ttfamily user\_read\_media / ``View media'' / Read; Subject: CurrentUser; Object: Media; requires Know(MediaToken)\\
\bottomrule
\end{tabular}
\caption{Worked example (MediaCMS). One field, \cd{friendly_token}, is traced from
observed traffic, to an \L1 node, to a semantic \cd{MediaToken}, to the \L2 ``View
media'' action---evidence to interface to meaning to business action.}
\label{fig:worked}
\end{figure}

\subsection{Agentic generation pipeline}
\begin{enumerate}[leftmargin=1.7em,itemsep=2pt]
\item \textbf{Freeze and deploy} (as above).
\item \textbf{Identify the loop and its edges} from source: the channel inbound path,
      the tool registry, and the policy chain---locating where context enters, exits,
      and originates.
\item \textbf{Extract records from source.} Map each agentic record (e.g.\ a terminal
      tool with \cd{effect_class: code_execution}, a sandbox \cd{execution_identity},
      and \cd{approval_required}).
\item \textbf{Classify trust and authority}: is the message \cd{trust_class:
      untrusted_data}; is the tool's \cd{effect_class} code execution; whose identity
      does it execute as.
\item \textbf{Draw the connections} (Figure~\ref{fig:ece}): wire
      \emph{enter}$\to$\emph{carry}$\to$\emph{exit} so every piece of content is
      traceable end to end.
\item \textbf{Runtime grounding.} Observe the live loop across entry, model call, tool
      dispatch, result, and exit while varying identity/requester, and confirm the
      runtime record matches the source records.
\item \textbf{Link to L1/L2.} Resolve each agentic identifier
      (\cd{invocation_surfaces[].l1_id}, \cd{model_config.writers[].l1_id},
      \cd{tool_bindings.l1_ids}, \cd{related_l2_ids}) to concrete \L1 nodes and \L2
      actions.
\item \textbf{Validate} (Section~\ref{sec:val}).
\end{enumerate}

\begin{figure}[ht]\centering
\begin{tikzpicture}[
  box/.style={draw, rounded corners, text width=2.9cm, minimum height=1.5cm,
              align=left, inner sep=5pt, font=\scriptsize},
  stage/.style={font=\scriptsize\bfseries, text=gold}]
\node[box, fill=agamber] (cs) {\textbf{context\_source}\\ Inbound channel message\\[2pt]\textit{trust: untrusted\_data}};
\node[box, fill=l1blue, right=1.0cm of cs] (ag) {\textbf{agent}\\ model decision\\[2pt]\textit{selects a tool}};
\node[draw, diamond, aspect=1.3, fill=green!8, right=0.6cm of ag, inner sep=1pt, font=\scriptsize] (ctl) {control};
\node[box, fill=redbox, right=0.6cm of ctl] (tl) {\textbf{tool: terminal}\\ effect: code\_execution\\ identity: sandbox\\ approval\_required};
\node[box, fill=l2green, right=1.0cm of tl] (os) {\textbf{output\_sink}\\ reply + exec result\\[2pt]\textit{audience: channel}};
\draw[-{Latex[length=2mm]}, navy, thick] (cs) -- (ag);
\draw[-{Latex[length=2mm]}, red!60!black, thick] (ag) -- (ctl);
\draw[-{Latex[length=2mm]}, red!60!black, thick] (ctl) -- (tl);
\draw[-{Latex[length=2mm]}, navy, thick] (tl) -- (os);
\node[stage, above=1pt of cs] {1 $\cdot$ ENTER};
\node[stage, above=1pt of ag] {2 $\cdot$ CARRY};
\node[stage, above=1pt of os] {3 $\cdot$ EXIT};
\end{tikzpicture}
\caption{Agentic lineage: an untrusted channel message enters as a context source,
the agent selects a tool, a control gates the edge, the terminal tool executes code as
the sandbox, and the result exits via an output sink. The root cause of agentic attack
vectors is authority on this path.}
\label{fig:ece}
\end{figure}

\subsection{Runtime grounding and its measurement}\label{sec:ground}
Both pipelines end by reconciling the authored records against a live, pinned
deployment. We make the \emph{strength} of that reconciliation explicit rather than
binary, through four instruments.

\paragraph{Evidence-depth ladder.} Each \L1 interface carries a grounding depth on a
fixed ladder---\cd{none} $\prec$ \cd{route-confirmed} (a request reached the handler,
e.g.\ a $4xx$ that proves the route exists) $\prec$ \cd{observed} (at least one $2xx$)
$\prec$ \cd{exercised} (observed across more than one identity or method). Each \L2
action carries a tier of \cd{interface} (its bound \L1 routes were observed) or
\cd{effect} (the action's state change was itself confirmed at runtime). The depth is a
property of the evidence, so a reader can see not only \emph{whether} a record is
grounded but \emph{how far}.

\paragraph{Completeness oracle.} To check that the authored surface is neither invented
nor badly under-counted, each \L1 set is diffed against the application's own route
list---either the served OpenAPI document (Oracle~A) or an in-container router dump
(Oracle~B). We report \emph{validity} (the share of checkable authored routes that are
real live routes) separately from \emph{coverage} (the share of the live surface that is
inventoried), and hold surface a static oracle structurally cannot enumerate (dynamically
mounted routers, packaged sub-APIs, static assets) out of the validity denominator so the
measure is not understated.

\paragraph{Multi-identity probing.} Interfaces are driven under more than one identity
(anonymous, user, admin, and application-specific roles such as Vocechat's
\cd{bot}/\cd{secret_holder}) and the per-identity status is recorded on the node, so an
authorization wall (e.g.\ $403$ as a member, $200$ as an admin) is captured as evidence
rather than inferred. Across the corpus \textbf{1{,}311} \L1 nodes carry status under two
or more identities.

\paragraph{Blocker taxonomy.} An agentic record that is not observed is tagged with
exactly one reason from a closed set: \cd{local_not_driven} (exercisable locally, not yet
driven), \cd{scheduling_required} (needs a timer/cron fire), \cd{external_integration}
(needs a real third-party credential a local stand-in cannot supply), or
\cd{crypto_blocked} (an auth/encryption payload that could not be reproduced). This turns
``not observed'' from a silence into a stated, checkable cause, and separates a
\emph{reachability} limit from a \emph{principled} one.

\subsection{Validation}\label{sec:val}
Validation enforces agreement automatically; a record that fails any check is not shipped.
\begin{itemize}[leftmargin=1.5em,itemsep=1pt]
\item \textbf{Traditional (L1/L2).} Closed vocabularies (predicates $\in$ the fixed set);
      semantic consistency (\cd{semantic_outputs} $=$ the set of mergeable resources);
      cross-layer closure (every \L2 \cd{l1_binding} resolves to a real \L1 route);
      uniqueness (no duplicate \L1 keys or action keys).
\item \textbf{Agentic.} Referential integrity (every \cd{*_id} points to a real record);
      lineage closure (every writer/output resolves); enum validity; existence of
      all referenced L1/L2 nodes; and evidence invariants (an observed record carries no
      blocker; every unobserved record carries exactly one valid blocker).
\end{itemize}
These checks run as committed scripts over the whole dataset, so an integrity regression
is caught rather than assumed; the current dataset passes with zero duplicate keys, zero
dangling cross-layer references, and all evidence invariants satisfied.

% ======================================================================
\section{Environments (Dataset)}
% ======================================================================

\begin{table}[ht]\centering\small
\caption{Traditional web environments and representative vulnerability classes.}
\label{tab:envweb}
\begin{tabularx}{\textwidth}{@{}l>{\raggedright\arraybackslash}X l@{}}
\toprule
\textbf{Environment} & \textbf{Tech stack} & \textbf{Vulnerability}\\
\midrule
VoceChat    & Rust REST (OpenAPI), SSE streams, Bot API (\cd{x-api-key}), admin webhooks & SSRF (CVE-2026-100893)\\
MediaCMS    & Django + React, REST, upload-driven media handling & RCE (CVE-2024-52004)\\
WriteFreely & Go templates, token REST, anonymous publishing, ActivityPub, OAuth & Local credential read (CVE-2025-24337)\\
Uptrace     & OTLP over gRPC/HTTP (project-DSN token), ClickHouse, Postgres & SQLi (CVE-2024-44905/44906)\\
Discourse   & Rails, Ember SPA, PSQL, Redis, Sidekiq, nginx & SSRF (CVE-2025-68662), IDOR (CVE-2026-26265), XSS (CVE-2025-48954)\\
Ghost       & Node, Ember admin, two APIs, theme upload, S3/GCS, webhooks & SQLi, token/2FA bypass, SSRF, RCE, path traversal\\
Rocket.Chat & Meteor/Node, REST + DDP WebSocket, Apps-Engine, SAML/CAS/OAuth & NoSQLi, unauthenticated file deletion\\
Keycloak    & Java/Quarkus, React console, admin REST, OIDC/SAML flows & Multiple authentication CVEs\\
Directus    & Node/Vue, REST + GraphQL, WebSocket, TUS uploads, URL import, Flows & Unauth.\ file upload/modification, SSRF, path traversal\\
Outline     & Node/TS + React, RPC-style REST, WebSocket, email-OTP, OAuth/API-key & OTP brute force, OAuth scope smuggling, IDOR, path traversal, XSS\\
\bottomrule
\end{tabularx}
\end{table}

\begin{table}[ht]\centering\small
\caption{LLM-integrated environments and representative vulnerability classes.}
\label{tab:envllm}
\begin{tabularx}{\textwidth}{@{}l>{\raggedright\arraybackslash}X l@{}}
\toprule
\textbf{Environment} & \textbf{Tech stack} & \textbf{Vulnerability}\\
\midrule
Open WebUI  & Python backend, Socket.IO, server-side Python, MCP tools, RAG & Multiple agentic tool vulnerabilities\\
LibreChat   & Node/Express + React, agents with permission/file models, MCP & Multiple MCP issues, unvalidated callback session\\
AnythingLLM & Node + React, HTTP API, agent WebSocket, developer/extension keys, workspaces, vector DBs, agent skills & SQL-agent injection; streaming-phase XSS$\to$RCE via LLM response injection; AgentFlows path traversal\\
LobeChat    & Next.js, tRPC, OIDC/Better Auth, PG/Redis/S3, agent \& MCP markets, URL/GitHub skill import, RAG & BOLA via chat-group agent operations; cross-user document disclosure via unscoped RAG; SSRF \& ReDoS in skill import\\
OpenClaw    & Node gateway (WS/HTTP/Control UI), chat channels + webhooks, exec tool (allowlist/approval), browser tool, skills registry, sub-agents, MCP + vector memory & Env-variable allowlist bypass; working directory injected into prompts; owner-only tools reachable via agent runs; prompt injection\\
\bottomrule
\end{tabularx}
\end{table}

% ======================================================================
\section{Results: Captured Attack Surface}
% ======================================================================

\begin{table}[ht]\centering\small
\caption{Captured traditional-web attack surface ($\Sigma$ = 1{,}397 endpoints,
980 semantic flows).}\label{tab:resweb}
\begin{tabular}{@{}lrr@{}}
\toprule
\textbf{Environment} & \textbf{L1 endpoints} & \textbf{L2 semantic flows}\\
\midrule
VoceChat    & 189 & 205\\
MediaCMS    & 110 & 104\\
WriteFreely & 101 &  30\\
Uptrace     &  54 &  36\\
Discourse   & 231 & 287\\
Ghost       & 112 &  59\\
Rocket.Chat & 154 &  46\\
Keycloak    & 125 &  50\\
Directus    & 115 &  68\\
Outline     & 206 &  95\\
\midrule
\textbf{Total} & \textbf{1{,}397} & \textbf{980}\\
\bottomrule
\end{tabular}
\end{table}

\begin{table}[ht]\centering\small
\caption{Captured LLM-integrated (agentic) attack surface, by record type
($\Sigma$ = 327 agentic records).}\label{tab:resllm}
\begin{tabular}{@{}lrrrrrrrr@{}}
\toprule
\textbf{Environment} & \textbf{Agents} & \textbf{Tools} & \textbf{Invoc.} &
\textbf{Context} & \textbf{Memory} & \textbf{Sinks} & \textbf{Edges} & \textbf{Controls}\\
\midrule
Open WebUI  & 4 & 22 & 5 & 16 & 7 &  8 & 13 & 12\\
LibreChat   & 3 &  9 & 3 & 12 & 5 &  9 &  9 &  7\\
AnythingLLM & 2 & 17 & 3 & 11 & 5 &  8 &  8 &  7\\
LobeChat    & 1 &  5 & 2 &  8 & 3 &  4 &  6 &  4\\
OpenClaw    & 3 & 27 & 7 & 12 & 5 & 10 & 12 & 13\\
\midrule
\textbf{Total} & \textbf{13} & \textbf{80} & \textbf{20} & \textbf{59} & \textbf{25}
& \textbf{39} & \textbf{48} & \textbf{43}\\
\bottomrule
\end{tabular}
\end{table}

\subsection{Runtime grounding and evidence depth}
The counts above are the \emph{size} of the dataset; this section reports its
\emph{evidence}. Table~\ref{tab:ground} grades the traditional-web corpus by the depth
ladder of Section~\ref{sec:ground}: every one of the 1{,}397 interfaces and 980 actions
is grounded (no ungrounded record), and the grading shows how far---358 interfaces
confirmed by an observed response ($2xx$), the remainder confirmed as reachable routes,
and 276 of the 980 actions confirmed by their observed effect rather than only through a
bound interface.

\begin{table}[ht]\centering\small
\caption{Traditional-web runtime grounding by evidence depth (all 1{,}397 \L1 / 980 \L2
records grounded).}\label{tab:ground}
\begin{tabular}{@{}lr@{\hspace{2.2em}}lr@{}}
\toprule
\textbf{\L1 interface depth} & \textbf{count} & \textbf{\L2 action tier} & \textbf{count}\\
\midrule
exercised ($2xx$, $\ge2$ identities/methods) &  55 & effect-observed    & 276\\
observed (at least one $2xx$)                 & 303 & interface-observed & 704\\
route-confirmed (handler reached)             & 1{,}039 & ---            & ---\\
not grounded                                  &   0 & not grounded       &   0\\
\midrule
\textbf{Total} & \textbf{1{,}397} & \textbf{Total} & \textbf{980}\\
\bottomrule
\end{tabular}
\end{table}

The agentic layer is harder to ground because all observation uses a \emph{local} model
with no cloud credentials, so some surface is unreachable by construction.
Table~\ref{tab:aground} reports, per application, how many of the gradable agentic
records\footnote{The 20 \texttt{invocation\_surfaces} are sub-fields of agent records and
are not independently graded, leaving 307 gradable records.} are runtime-observed, and
accounts for every unobserved record by its blocker.

\begin{table}[ht]\centering\small
\caption{Agentic runtime grounding: observed records per application, and the blocker
accounting for the 178 unobserved.}\label{tab:aground}
\begin{tabular}{@{}lr@{\hspace{2.6em}}lr@{}}
\toprule
\textbf{Application} & \textbf{observed\,/\,total} & \textbf{Unobserved by blocker} & \textbf{count}\\
\midrule
AnythingLLM & 47 / 58 & \cd{local_not_driven} (exercisable, not yet driven) & 125\\
Open WebUI  & 36 / 82 & \cd{crypto_blocked} (irreproducible auth payload)   &  31\\
OpenClaw    & 22 / 82 & \cd{external_integration} (needs real OAuth)        &  15\\
LibreChat   & 24 / 54 & \cd{scheduling_required} (needs a timer fire)       &   7\\
LobeChat    &  0 / 31 & & \\
\midrule
\textbf{Total} & \textbf{129 / 307} & \textbf{Total unobserved} & \textbf{178}\\
\bottomrule
\end{tabular}
\end{table}

Separating the blockers distinguishes a \emph{principled} ceiling from a
\emph{reachability} limit. Only 38 records are a principled ceiling---the 31
\cd{crypto_blocked} (LobeChat's encrypted auth payload could not be reproduced) and the 7
of the external integrations that require a genuine third-party OAuth identity a local
stand-in cannot mint. The remaining 140 are reachability-limited: they are observable in
principle but were not driven here because they depend on the reliability of a local
model emitting a native tool call, or on activating a dormant service. Where a local
stand-in \emph{can} substitute for an external service without a secret (image
generation, a Model Context Protocol server, a scheduler, or a direct create/read of a
store), those records were driven to an observed effect; image generation and MCP are
therefore no longer categorical blockers.

\paragraph{Completeness.} Table~\ref{tab:complete} reports, for the applications whose
own route list could be dumped, the validity and coverage of the authored \L1 set
(Section~\ref{sec:ground}). High validity confirms the authored routes are real; coverage
is reported against the whole live surface and is deliberately partial where the dataset
scopes to a subset of a very large API (e.g.\ Discourse, Open WebUI), which the table
marks.

\begin{table}[ht]\centering\small
\caption{Authored-surface completeness against each application's own route oracle
(A = served OpenAPI, B = in-container router dump).}\label{tab:complete}
\begin{tabular}{@{}lccrr@{}}
\toprule
\textbf{Application} & \textbf{Oracle} & \textbf{live routes} & \textbf{validity} & \textbf{coverage}\\
\midrule
AnythingLLM & B &   67 & 100.0\% & 100.0\%\\
VoceChat    & A &  159 &  98.7\% &  98.7\%\\
Outline     & B &  213 & 100.0\% &  93.9\%\\
Directus    & A &   65 &  92.5\% &  95.4\%\\
Ghost       & B &  203 & 100.0\% &  43.8\%$^\ast$\\
Discourse   & B & 1{,}732 &  99.6\% &  13.3\%$^\ast$\\
Open WebUI  & B &  587 & 100.0\% &   2.9\%$^\ast$\\
\bottomrule
\end{tabular}

{\footnotesize $^\ast$ Dataset is deliberately scoped to a subset of the live API, so low
coverage is by design; for Ghost the oracle also surfaced a genuine gap (admin-API
families not yet inventoried).}
\end{table}

% ======================================================================
\section{Discussion and Key Insights}
% ======================================================================

\subsection{Traditional attack surface}
\begin{enumerate}[leftmargin=1.6em,itemsep=2pt]
\item \textbf{The exhaustive ground truth exposes surfaces crawlers miss}---notably
      non-HTTP entries (WebSocket/DDP, gRPC) and endpoints gated behind state or a
      workflow. These do not appear in a browser-centric crawl.
\item \textbf{The same business action is exposed through different surfaces}---an
      internal UI API, a developer/API-key API, and an embed widget can all realise one
      action. This is exactly where an automated scanner under-counts: it discovers one
      surface and stops. The schema binds all three to a single \L2 action, making the
      duplication explicit.
\end{enumerate}

\subsection{Agentic attack surface}
\begin{enumerate}[leftmargin=1.6em,itemsep=2pt]
\item \textbf{The root cause of attack vectors is authority}: untrusted data reaches a
      model decision, which reaches a privileged tool. By treating the execution edges
      as a traversable path, the schema lets an investigator enumerate these vectors
      directly rather than inferring them.
\item \textbf{Controls must be enforced, not merely present.} For each
      untrusted$\to$privileged path, the schema records whether the controls on that
      path (\cd{edge.control_ids} / \cd{tool.control_ids}) are actually enforced---turning
      a qualitative safety question into a checkable property.
\end{enumerate}

% ======================================================================
\section{Future Work}
% ======================================================================
\begin{enumerate}[leftmargin=1.6em,itemsep=2pt]
\item \textbf{Benchmark agents.} Define discovery metrics---precision/recall of surface
      discovery, accuracy of semantic descriptions, and correctness of cross-layer
      relationships---and run LLM crawlers and agentic scanners against all fifteen
      applications, scored against this ground truth.
\item \textbf{Extend schema coverage} to transports modelled only partially here:
      WebSocket/DDP, gRPC, and GraphQL.
\item \textbf{Close the reachability-limited grounding gap.} Runtime grounding is now
      established across the corpus (Table~\ref{tab:ground}); the residual is the 140
      reachability-limited agentic records---driving the model-dependent tool calls with
      a more reliable local model, activating the dormant services, and standing up the
      remaining secret-free local stand-ins---and the small set of source-first LLM web
      surfaces not yet driven. The 38 principled-ceiling records (irreproducible crypto
      and genuine third-party OAuth) are documented as out of scope for a local,
      credential-free environment.
\end{enumerate}

% ======================================================================
\section{Conclusion}
% ======================================================================
This project delivers an evidence-backed, multi-application ground-truth dataset
(1{,}397 endpoints, 980 semantic flows, 327 agentic records) and a repeatable
annotation-and-review workflow that reconciles source and runtime evidence, grades each
record by observation depth, and checks every cross-layer reference. It enables rigorous benchmarking of security agents'
ability to \emph{generate} an attack surface, provides a concrete framework for
measuring agentic risk, and offers a reusable methodology and dataset for future
extensions.

% ======================================================================
\begin{thebibliography}{9}
\bibitem{owaspapi} OWASP Foundation. \emph{OWASP API Security Top 10 (2023)} ---
  API1:2023 Broken Object Level Authorization.
\bibitem{wstg} OWASP Foundation. \emph{Web Security Testing Guide (WSTG)} ---
  authentication and authorization testing.
\bibitem{owaspllm} OWASP Foundation. \emph{OWASP Top 10 for Large Language Model
  Applications} (2025).
\bibitem{atlas} MITRE. \emph{ATLAS: Adversarial Threat Landscape for Artificial
  Intelligence Systems}.
\bibitem{strips} R.~E. Fikes and N.~J. Nilsson. STRIPS: A New Approach to the
  Application of Theorem Proving to Problem Solving. \emph{Artificial Intelligence},
  1971.
\bibitem{pddl} D.~McDermott et al. \emph{PDDL --- The Planning Domain Definition
  Language}. 1998.
\bibitem{openapi} OpenAPI Initiative. \emph{OpenAPI Specification}.
\end{thebibliography}

\end{document}
